\project{18}{Safe states, and the workspace sensor that triggers one}{A way to stop}{Safe-state machine, latching, entry actions, a presence sensor as a trigger}

\begin{keyfacts}
\item[Adds to the node] A way to stop. The node gains a latched safe state
reachable from everywhere, three independent triggers, and a written answer to
the question a joint has to answer before it goes near anybody
\item[Peripherals] The ranging shield as the workspace trigger, both on-die
watchdogs, and one output that asserts the safe state without software
\item[Depends on] Chapter 5 for position, chapter 7 for the output stage,
chapter 12 for bus loss, chapter 16 for following error
\item[Real or modelled] \textbf{Real trigger, modelled reaction.} The sensor,
the latch, the watchdogs and the output are real. \textbf{There is no motor and
no brake}, so the reaction is asserted and logged rather than felt, and the
chapter never implies otherwise
\item[Difficulty] 4 of 5
\item[Effort] Three evenings, one of which is reading standards and finding out
what may and may not be claimed from them
\item[Deliverable] Two safe states rather than one, a latch that meets a
published clause by clause list, a trigger with a measured reaction time, and a
document that says exactly what this node is and is not
\end{keyfacts}

\subsection*{Why this chapter}

A joint that can move is a joint that has to be able to stop, and the word for
what it stops into has a definition. That definition is the first surprise of
this chapter, because it contains no engineering content at all.

The definitions part of the general functional safety standard gives it in
seven words: the state of the equipment under control when safety is achieved.
It does not say power removed. It does not say brake applied. \textbf{What
counts as the safe state is an output of the risk assessment, not of the
firmware.} And the note attached to the definition is the load-bearing sentence
for a robot joint: for some situations a safe state exists only so long as the
equipment is continuously controlled.

\begin{note}{A joint has two safe states and they are not interchangeable}
For a horizontal joint, or one whose gearing will not back-drive, removing
torque is a safe state and it persists with no energy at all. \textbf{For a
vertical joint carrying a payload, removing torque is not a safe state}, because
the axis then descends. There the safe state is the brake engaged with torque
removed, and it is reached through a controlled deceleration rather than
instantly. Torque removed and brake held are two different safe states, the
drive standard names them and the definitions standard names neither, and a
design that has only one of them has only solved half of the problem. This
chapter builds the machine for both and says which one this bench can
demonstrate.
\end{note}

\subsection*{Prior art and what to reuse}

\begin{table}[H]\centering\small
\begin{tabular}{p{34mm}p{46mm}p{38mm}p{20mm}}
\toprule
Source & What it gives & What it does not & Licence \\
\midrule
The definitions part of the general standard, edition 2.0, Friday 30 April 2010 & The safe
state definition itself, \textbf{read from the official preview at printed page
11}, with the note about continuous control that this chapter is built around.
Also the tool classification into three types, read directly & Paywalled beyond
the preview. It says nothing about how to implement anything, deliberately &
paid \\
The emergency stop standard, third edition, 2015 & \textbf{Clause 4.1.1, read
directly}: six requirements the latch in this chapter is written against, and
clause 4.1.1.3, which says plainly that it is a complementary measure and does
not substitute for safeguarding & \textbf{Its clause on stop categories sits
beyond the end of the free preview}, so that part is cited by clause number and
attributed. \textbf{Its own normative references point at an edition of the
electrical equipment standard that is now two editions behind} & paid \\
One silicon vendor's free industrial functional safety guide & The distinction
firmware engineers most often miss, in a per-device table: \textbf{random
hardware capability against systematic capability}, and the architecture
arithmetic in a line each. Then the list of mechanisms that are actually
firmware & It is about a different vendor's parts, so the arithmetic transfers
and the device table does not & free, no registration \\
The performance level calculator from an institute & Models the control structure
and computes the attained level, free of charge, version 3.0.4 of Friday 31
October 2025, supporting the 2023 edition & \textbf{Free to use and to pass on,
but modification and re-hosting are not permitted.} Link it, never vendor it &
free, restricted \\
The space agency's software assurance handbook, topic 8.5 & An eight step
analysis method, \textbf{a catalogue of software failure modes that reads as if
written for a bus-connected joint node}, criticality categories, and four
downloadable worksheet templates & Nothing certifiable. \textbf{It is a United
States government work with no copyright notice, which makes it the safest
material in this volume to adapt into a public repository} & public domain \\
A permissively licensed bus stack & A working implementation of the safety
data object wire format: a periodic broadcast of two frames, the second carrying
the data bit-wise inverted on a different identifier & \textbf{It claims coding
standard conformance with documented exceptions, which is not a functional safety
certification}, and it makes no certification claim anywhere & Apache-2.0 \\
A maintained hierarchical state machine library & Entry, run and exit actions,
history, timers, parallel states, a kernel dispatcher, and coding standard checks
in its own continuous integration & C++ where this node is C, which is a real
mismatch to state rather than to paper over & MIT \\
\bottomrule
\end{tabular}
\caption{Prior art for chapter 18. The fifth row is the most useful thing in the table for a reader building a portfolio: a complete, free, quotable analysis method with worksheets, from a source whose material carries no copyright restriction at all.}
\end{table}

\subsection*{What the node gains}

A latched safe state that is reachable from every other state, three independent
ways into it, an entry action that does not depend on the scheduler still
running, and a written statement of what this node is. That last item is the one
that matters most in a portfolio: a node that says clearly that it implements a
safe-state machine and makes no integrity claim is more credible than one that
implies a claim it cannot support.

\subsection*{Parts from the inventory}

\begin{table}[H]\centering\small
\begin{tabular}{p{40mm}p{66mm}p{40mm}}
\toprule
Part & Role & Interface \\
\midrule
X-NUCLEO-53L8A1, 8 by 8 ranging & \textbf{The workspace trigger, and it is
real.} Multizone ranging means a region rather than a point, which is what a
workspace is & The shield header \\
NUCLEO-H7A3ZI-Q & The machine, the latch and the entry action & Its own outputs \\
The independent watchdog & Clocked from the low speed internal oscillator, so it
survives a stopped main clock & On-die \\
The window watchdog & Faults on a refresh that is \textbf{too early} as well as
too late, which catches a loop running at the wrong rate & On-die \\
An external watchdog & \textbf{Not on the bench, and named as absent.} Both
on-die watchdogs share the chip's fate, so a genuinely independent one is
external & n/a \\
A brake, and a motor & \textbf{Neither exists here.} The reaction is asserted on
a pin and logged, and no figure in this chapter implies a mechanism moved & n/a \\
\bottomrule
\end{tabular}
\caption{Inventory items used in chapter 18. Two rows are absences and both are load-bearing. The missing external watchdog limits what may be claimed about diagnostic independence; the missing brake means the second of this chapter's two safe states is designed, tested in logic, and never demonstrated physically.}
\end{table}

\subsection*{System architecture}

\diagram{j18_arch}{The safe-state machine, its three independent triggers and the
latch. Two rules shape it, and neither has a citable source, so both are
presented as engineering practice rather than dressed up as requirements: the
safe state is reachable from every other state, and once entered it is held
until an intentional reset. The published clause list that the reset is written
against is in the figure beside it.}

\subsection*{Peripheral configuration}

\begin{table}[H]\centering\small
\begin{tabular}{p{22mm}p{26mm}p{24mm}p{30mm}p{30mm}}
\toprule
Peripheral & Mode & Clock source & Pins and function & Interrupt and transfers \\
\midrule
Ranging sensor & Multizone, continuous, with its own interrupt & Its own & The
shield's bus pins plus one interrupt line & Interrupt on a new frame, never
polled \\
Independent watchdog & Enabled before anything else, and never disabled & The low
speed internal oscillator & None & Survives a stopped main clock \\
Window watchdog & Refreshed once per control period, inside the window & From
the bus clock & None & \textbf{Faults on an early refresh too} \\
Output stage & \textbf{Break input already proven in chapter 7} & Unchanged &
Unchanged & \textbf{Asserts with interrupts disabled}, which chapter 7 measured \\
Safe-state output & One pin, asserted by the same hardware path & n/a & One pin,
plus its indicator & No software in the assertion path \\
\bottomrule
\end{tabular}
\caption{Peripheral configuration for chapter 18. The fourth row is chapter 7's work being spent: the break input was proven there to remove the outputs with interrupts disabled, which means the entry action into the safe state does not depend on the scheduler, on a task, or on anything continuing to run.}
\end{table}

\subsection*{Wiring}

\diagram{j18_wiring}{What is on the bench, and what a real safety chain has that
this one does not. The trigger, the latch and the output are real and single
channel. A chain that could support an integrity claim is dual channel with
diagnostics, and both the second channel and the external watchdog are drawn as
absent rather than implied. The figure is the honest picture, and it is also the
reason the chapter's own claim is carefully worded.}

\subsection*{Memory and timing budget}

\begin{table}[H]\centering\small
\begin{tabular}{p{44mm}p{28mm}p{28mm}p{30mm}}
\toprule
Quantity & Budget & Measured & Margin \\
\midrule
Flash, this chapter & 7 kB & not measured & not measured \\
Static memory, this chapter & 2 kB, mostly the ranging frame & not measured & not measured \\
Trigger to output asserted & under 2 ms & not measured & not measured \\
Ranging frame period & 15 Hz, which bounds the first term above & not measured & not measured \\
Controlled deceleration, in logic & 200 ms, then torque removed & not measured & not measured \\
Latch, once set & \textbf{forever, until an intentional reset} & n/a & n/a \\
Reset to running & one full self-check, never immediate & not measured & not measured \\
\bottomrule
\end{tabular}
\caption{The budget for chapter 18. The third row is the whole reaction time and most of it is the sensor: a trigger cannot be acted on before it arrives, and a 15 hertz frame rate puts a floor under the reaction that no amount of firmware speed removes. Stating that is more useful than reporting the microseconds the software takes.}
\end{table}

\subsection*{Firmware design (UML)}

\diagram{j18_uml}{Two safe states, and what each one costs to reach. Removing
torque is immediate and persists with no energy, and it is not a safe state at
all for an axis that would descend. Brake held is reached through a controlled
deceleration with power still available, and only then is power removed. The
three columns on the right are how those map onto the published stop categories
and the drive standard's sub-functions.}

Three rules, and the first two have no citable source.

\textbf{The safe state is reachable from every other state.} No state may be a
place from which the safe state cannot be entered. This is real and widely
practised, and this volume could find no published source stating it, so it is
presented as engineering practice with the standards cited for the surrounding
requirements and never attributed to a document nobody has read.

\textbf{A fault latches.} The condition that caused entry is recorded, the state
is held, and neither clears because the condition went away. The same applies:
practised everywhere, named in no source this volume could obtain.

\textbf{The entry action does not depend on software continuing to run.} Chapter
7 proved the break input removes the outputs with interrupts disabled. That
proof is what makes the entry action trustworthy, and a design whose safe state
is entered by a task is a design whose safe state depends on the scheduler.

\subsection*{Data flow (ASCII)}

\begin{asciiart}
  three independent triggers, and none of them can be disabled by the others
  ---------------------------------------------------------------------
        |                      |                       |
  workspace violated     following error         bus loss or bus-off
  (the ranging shield,   outside tolerance       (chapters 12 and 13)
   15 Hz, real)          (chapter 16)
        |                      |                       |
        +----------------------+-----------------------+
                               |
                               v
                      ENTER THE SAFE STATE
                               |
        +----------------------+-----------------------+
        |                                              |
   torque removed                              controlled deceleration
   immediate, and it PERSISTS                  power still available
   with no energy                              then power removed, and
        |                                      THE BRAKE HOLDS
        |                                              |
   for a horizontal or non-backdrivable          for a vertical axis
   joint, this IS the safe state                 carrying a payload,
        |                                        the one above is NOT
        v                                              v
                      LATCHED, and recorded with its cause
                               |
                               v
             reset only by an intentional human action, and
             THE RESET ITSELF DOES NOT START ANYTHING
                               |
                               v
                     a full self-check, then running

  on this bench: the trigger is real, the latch is real, the output pin is
  real, and there is no motor and no brake, so the reaction is asserted and
  logged rather than felt
\end{asciiart}

\subsection*{Repository layout}

\begin{plaincode}
joint-node/
  src/
    safety/
      safe_state.c safe_state.h   # + the machine, the latch, the entry action
      triggers.c                  # + three, independent, none maskable
      workspace.c                 # + the ranging shield as a region, not a point
      selftest.c                  # + what must pass before a reset takes effect
      wdg.c                       # + both on-die watchdogs, and a liveness mask
    control/ sense/ act/ bus/ mw/ estimate/ time/ node/ bsp/  update/
  doc/
    safe-states.md                # + BOTH of them, and which one applies when
    what-this-node-is-not.md      # + the claim, carefully worded
    fmea/                         # + the agency's worksheets, filled in
  test/
    test_safe_state.c             # + reachability from every state, on the host
    test_latch.c                  # + it does not clear when the cause clears
  host/  third_party/  tools/  proto/  README.md
\end{plaincode}

\subsection*{Steps}

\begin{steps}
\item \textbf{Write down which safe state applies, before writing the machine.}
This is a risk assessment output and firmware cannot decide it. Two lines of
honesty here save a design.
\begin{plaincode}
doc/safe-states.md

  horizontal or non-backdrivable axis
      safe state      torque removed
      persists        yes, with no energy at all
      reached         immediately
  vertical axis carrying a payload
      safe state      BRAKE ENGAGED with torque removed
      persists        as long as the brake holds
      reached         controlled deceleration FIRST, then power removed
      why             removing torque alone is not a safe state here: the
                      axis descends
  on this bench       the first, demonstrated. The second, designed and
                      tested in logic only, because there is no brake
\end{plaincode}

\item \textbf{Keep the three axes apart, because most requirements that arrive
have them tangled.} They answer different questions and they come from different
documents.
\begin{plaincode}
axis                what it specifies                        values
stop CATEGORY       what the stopping action does to
                    actuator power                           0, 1, 2
performance LEVEL   how RELIABLY the function is performed   a to e
integrity LEVEL     the same question, other document        1 to 3 (4)

a well formed requirement names ONE from the first and ONE from the others:
  "emergency stop shall be stop category 1, achieved with PL d"

THERE IS NO STOP CATEGORY 3 AND NO PL 1.
A category 0 stop implemented with a single undiagnosed contactor is still
a category 0 stop, at a poor performance level. The two are independent.
\end{plaincode}
\textbf{This volume prints no cross-map between the reliability scales.} Such
tables circulate widely, they are a convenience rather than a normative
equivalence, and no primary verified mapping was obtained during this volume's
research. Printing one would be repeating something rather than knowing it.

\item \textbf{Get the mapping to the drive standard right, because it is the
cleanest single fact in this area.} Stop category 0 corresponds to safe torque
off, category 1 to safe stop 1, and category 2 to safe stop 2. The drive
standard's second edition replaced the phrase safety function by \textbf{safety
sub-function} throughout, which is why drive documentation reads oddly to
somebody who learned the older wording.

\item \textbf{Write the latch against the published clause, line by line.} The
emergency stop standard's clause 4.1.1 gives six requirements and they are worth
copying into the code as comments.
\begin{ccode}
/* safe_state.c: each line below is a requirement, not a preference. */
/* 1. initiated by a single human action                                */
/* 2. available and operational at all times, in every mode             */
/* 3. overrides all other functions without impairing other protections */
/* 4. maintained until manually reset                                   */
/* 5. no start command is effective on the stopped operations           */
/* 6. reset by intentional action, and THE RESET SHALL NOT RESTART      */
void safe_state_enter(trigger_t why)
{
    act_break_assert();            /* chapter 7: with interrupts disabled */
    g_safe.latched = true;
    g_safe.cause   = why;          /* recorded, and it survives the state */
    g_safe.count++;                /* lifetime, and it survives a reset   */
}
\end{ccode}
The same clause adds, in 4.1.1.3, that this is a \textbf{complementary
protective measure} and does not substitute for safeguarding or for other safety
functions. That sentence belongs in the documentation of any node that
implements one, because it is the sentence that stops a stop function from being
treated as the whole safety case.

\item \textbf{Make the workspace sensor a region rather than a point.} An eight
by eight ranging frame is sixty-four distances, and the useful question is not
what any one of them says.
\begin{ccode}
/* workspace.c: a zone policy with hysteresis, and a minimum valid count. */
unsigned close = 0, valid = 0;
for (unsigned z = 0; z < 64; z++) {
    if (frame.status[z] != RANGE_VALID) continue;
    valid++;
    if (frame.mm[z] < g_ws.enter_mm) close++;
}
if (valid < g_ws.min_valid) { workspace_unknown(); return; }   /* not safe */
if (close >= g_ws.enter_zones) workspace_violated();
else if (close == 0 && g_ws.state == VIOLATED &&
         all_beyond(&frame, g_ws.exit_mm)) workspace_clear();  /* hysteresis */
\end{ccode}
\textbf{Too few valid zones is not the same as nothing being there.} A sensor
that cannot see is a sensor whose output must not be read as all clear, and that
distinction is one line of code and the difference between a safety function and
a decoration.

\item \textbf{Prove reachability on the host, from every state.} This is the
rule with no source behind it, and a test is the only thing that keeps it true
as the machine grows.
\begin{ccode}
/* test_safe_state.c: for every state, the safe state must be reachable. */
for (state_t s = 0; s < STATE_COUNT; s++) {
    machine_force_state(s);
    safe_state_enter(TRIGGER_TEST);
    TEST_ASSERT_EQUAL(STATE_SAFE, machine_state());
}
\end{ccode}

\item \textbf{Prove the latch does not clear when the cause clears.} This is the
second unsourced rule and the second one-line test.
\begin{shellcode}
python host/provoke.py --workspace-violate --hold 0.5 --then-clear
# expected: safe state entered at t0, and STILL entered at t0 + 30 s
# expected: cause reads WORKSPACE, not NONE, after the sensor reads clear
\end{shellcode}

\item \textbf{Use both watchdogs, for different reasons, and say what neither of
them proves.} The independent one is clocked from the low speed oscillator and
survives a stopped main clock. The window one faults on a refresh that arrives
too early, which catches a loop running faster than it should rather than one
that has stopped.
\begin{ccode}
/* wdg.c: a liveness mask, because refreshing from a timer proves nothing. */
void wdg_service_from_supervisor(void)
{
    if (g_live.mask == LIVE_ALL_TASKS) {    /* every task checked in */
        wdg_refresh();                      /* inside the window, not early */
        g_live.mask = 0;
    }                                       /* otherwise: let it fault */
}
\end{ccode}
\textbf{Both watchdogs are on-die and share the chip's fate.} A genuinely
independent watchdog for a safety argument is an external component, this bench
does not have one, and the chapter says so rather than letting two on-chip
timers stand in for independence.

\item \textbf{Read the vendor's safety package claim correctly, and do not
repeat the garbled version.} The architecture claim is specific and it is
frequently misquoted.
\begin{plaincode}
what the vendor's page says
  "SIL2 safety functions can be implemented with a single MCU;
   SIL3 safety functions implementation requires two MCUs in a 1oo2 scheme."
what several write-ups say instead
  the same sentence with the redundancy scheme changed, which reverses the
  meaning of the second half
what the package contains
  a safety manual, a qualitative analysis, a static failure rate report,
  and a self test library delivered AS OBJECT CODE
what could NOT be verified for this volume
  the licence text, the distribution mechanism, the current version, and
  WHETHER THIS EXACT PART IS COVERED AT ALL
\end{plaincode}
That last line is why nothing from that package is used here. An object-code
library bound by conditions of use is not something a book can tell a reader to
drop into a repository they intend to publish.

\item \textbf{Be equally careful about the kernel.} The kernel this volume uses
is permissively licensed and its own documentation makes no safety certification
claim. A separate, certified product shares a functional foundation with it and
is a \textbf{completely redesigned product} rather than a certified edition of
the same code. Writing the certified product's name next to the free one, or
calling either a certified version of the other, is the single most common
error in this area.

\item \textbf{Do the failure analysis with free material and a spreadsheet.} The
open tooling for this is immature and none of it is worth the dependency.
\begin{plaincode}
method        the space agency's handbook, topic 8.5: eight steps, and it is
              explicitly based on the cancelled military standard that
              everything else descends from
worksheets    four templates, downloadable, no copyright notice, United
              States government work: adapt them and attribute
failure modes its catalogue reads as if written for this node: out of range
              values, missing inputs, overwritten memory, data collisions,
              command omission, incorrect sequence, illegal commands,
              timing issues, and hardware and software interaction failures
tooling       none. A spreadsheet under version control beats every open
              tool this volume could find, and two of those carry six or
              seven commits in total
\end{plaincode}
One contrast is worth knowing about before somebody argues with you: the
international standard keeps the risk priority number and adds an alternative
beside it, while the automotive sector handbook is widely reported to have
dropped it for a lookup table. \textbf{Two engineers working to different sector
standards are told opposite things about the same number}, and neither is wrong.

\item \textbf{Write down what this node is not.} One page, and it is the
deliverable that makes the rest credible.
\begin{plaincode}
doc/what-this-node-is-not.md

  it IS        a node with a latched safe state, three independent
               triggers, an entry action proven to run with interrupts
               disabled, and a documented reaction time
  it is NOT    a safety-rated device. Single channel, no diagnostic
               coverage figure, no external watchdog, no assessment, no
               certificate, and no integrity level claimed
  the gap      dual channel with fault tolerance 1 reaches the higher
               levels; single channel with fault tolerance 0 does not, and
               that is architecture rather than effort
  and note     random hardware capability and systematic capability are
               different things. Systematic capability is about the
               development process; random hardware capability is about
               diagnostic coverage and failure rates. A part may be rated
               differently on each, and most are
\end{plaincode}
\end{steps}

\diagram{j18_timing}{Two stop categories drawn in time, and the latch that
follows both. Above, power removed immediately, which is a safe state for one
mounting and the opposite of one for another. Below, a controlled deceleration
with power still available, then power removed once stopped, which is what a
gravity-loaded axis needs. The bar at the bottom is the latch: it outlasts its
cause by design, and the reset that ends it starts nothing.}

\subsection*{Build, flash and debug}

\diagram{j18_data}{Three scales that get tangled in almost every requirement
document, kept apart, with what each one comes from and what values it actually
has. Below them, this volume's evidence marking for the clauses in this chapter:
what was read from a primary document, what was attested by agreeing secondary
sources, and what is presented as engineering practice with no source at all.}

\begin{shellcode}
cmake --build build -j && ctest --test-dir build/host
probe-rs run --chip STM32H7A3ZITx build/firmware.elf
python host/provoke.py --workspace-violate --measure-reaction --runs 200
python host/provoke.py --sensor-blind --expect workspace_unknown
\end{shellcode}

\begin{note}{When a stop function looks like a safety case}
It is not one, and the standard says so in a clause of its own: an emergency
stop is a complementary protective measure and does not substitute for
safeguarding or for other safety functions. A node with a good stop function and
no assessment has a good stop function. The temptation in a portfolio is to let
the reader infer more than that, and the correction is a single page saying
plainly what the node is not, which costs nothing and is the most credible thing
in the repository.
\end{note}

\subsection*{Verification and acceptance criteria}

\begin{itemize}
\item Both safe states are written down with the condition under which each
applies, and the bench says which one it can demonstrate.
\item The safe state is reachable from every other state, proven by a host test
that iterates over the state list rather than over a remembered subset.
\item The latch does not clear when its cause clears, proven by a thirty second
hold after the sensor reads clear.
\item The cause is recorded, and a lifetime count survives a reset.
\item The entry action asserts the output with interrupts disabled, which
chapter 7 already measured.
\item Three triggers are independent and none can be masked by the others.
\item A sensor that cannot see produces an unknown state, not a clear one,
proven by covering it.
\item Reaction time from trigger to output is measured over two hundred runs and
reported as median, 99.9th percentile and maximum.
\item The reset requires an intentional action, runs a full self-check, and does
not itself start anything.
\item The window watchdog is refreshed from a liveness mask rather than from a
timer, and a deliberately stalled task causes a fault.
\item The document saying what the node is not exists, and names single channel,
the absent external watchdog and the absent assessment.
\item No cross-map between the reliability scales is printed anywhere in the
repository.
\end{itemize}

\subsection*{Variants}

\begin{table}[H]\centering\small
\begin{tabular}{p{24mm}p{26mm}p{44mm}p{22mm}p{20mm}}
\toprule
Axis & Variant & What changes & Cost & Built in full in \\
\midrule
Safe state & Torque removed & Immediate, and persists with no energy & \textbf{Not a safe state at all for an axis that would descend} & Here, demonstrated \\
Safe state & Brake held after a controlled deceleration & The only correct answer for a gravity-loaded axis & A brake, its release circuit, and a deceleration that must complete & Here, in logic only \\
Trigger & Workspace sensor & A region, with hysteresis and a blind-sensor state & 15 Hz, which is most of the reaction time & Here \\
Trigger & Following error & Free: chapter 16 already computes it & Choosing a tolerance that is neither deaf nor twitchy & Here \\
Trigger & Bus loss & Free: chapter 12 already detects it & A policy decision about how long is too long & Here \\
Watchdog & From a timer & One line, and it proves the scheduler runs & \textbf{It refreshes happily while a task is stalled} & Nowhere \\
Watchdog & From a liveness mask & Every task checks in, or the fault happens & A bit per task and one supervisor & Here \\
Watchdog & External & The only kind that does not share the chip's fate & A component this bench does not have & Nowhere, and named \\
Architecture & Single channel & What this bench is & Reaches the lower integrity levels at best, by architecture & Here \\
Architecture & Dual channel with diagnostics & What an integrity claim needs & A second channel, comparison, and an assessment & Nowhere \\
Bus safety & A safety layer over the bus & Two frames, the second bit-wise inverted, so small parts can participate & \textbf{The published layer is for classic frames, and no equivalent for the flexible-data format was established} & Nowhere, and the gap is named \\
State machine & Hand written, flat & What this node has: small, auditable, testable & It will not scale to a hierarchy & Here \\
State machine & A library & History, parallel states, a kernel dispatcher, coding standard checks & C++ where this node is C & Nowhere, and the mismatch is stated \\
\bottomrule
\end{tabular}
\caption{Variants for chapter 18. The two safe-state rows at the top are the chapter in miniature: the same design decision is correct in one mounting orientation and incorrect in another, and nothing in the firmware can tell which one it is in.}
\end{table}

\subsection*{Pitfalls}

\begin{itemize}
\item Assuming the safe state is power removed. The definition says nothing of
the kind, and for a gravity-loaded axis power removed is the opposite of safe.
\item Having one safe state where a joint needs two.
\item Tangling the stop category with a reliability scale. They answer different
questions and come from different documents.
\item Printing a cross-map between the reliability scales as though it were
normative.
\item Writing stop category 3, or performance level 1. Neither exists.
\item Refreshing a watchdog from a timer, which proves only that the scheduler
runs while a task is stalled.
\item Treating two on-die watchdogs as independent. They share the chip's fate.
\item Reading too few valid ranging zones as all clear.
\item Letting a reset restart the machine. The standard says the reset shall not
initiate a restart, in the same clause that requires the reset.
\item Repeating the garbled version of a vendor's architecture claim, which
circulates widely and reverses the meaning of its second half.
\item Calling a certified kernel product a certified version of the free one.
They share a functional foundation and are separate, redesigned products.
\item Implying an integrity claim by silence. The cure is one page saying what
the node is not.
\item Citing the robot safety series by its old title. It is now titled
Robotics, and getting that wrong dates a document at a glance.
\item Writing that the collaborative robot technical specification is withdrawn.
It is published, was confirmed in 2022, and remains current until its replacement
appears. Several secondary sources say otherwise and they are wrong.
\item Citing the electrical equipment standard or the machinery functional
safety standard without their amendments. A bare citation is incomplete in 2026.
\item Citing the coding guidelines as the 2012 edition with amendments. A
consolidated edition superseded that arrangement, and a further edition is
current as of \pubdate{March 2025}.
\end{itemize}

\subsection*{Best practices applied}

\begin{itemize}
\item A definition is read from its own document rather than from what everybody
knows it says.
\item A design decision that firmware cannot make is identified as such and sent
back to the risk assessment.
\item Two practices with no citable source are named as engineering practice
rather than attributed to documents nobody has read.
\item Each claim carries its evidence class: read directly, attested by agreeing
sources, or practice.
\item An entry action is made independent of the scheduler and that independence
is proven, not assumed.
\item A sensor's inability to see is a distinct state from its seeing nothing.
\item A vendor claim is quoted exactly, and the widely circulating garbled
version is named.
\item What the node is not is written down, which is what makes what it is
believable.
\item Free, unrestricted public material is preferred over paid material
wherever it is adequate, and the analysis method chosen here is both.
\end{itemize}

\subsection*{Stretch goals}

\begin{itemize}
\item Add an external watchdog and measure what it catches that the on-die pair
does not. It is a small component and it changes what may be claimed.
\item Implement the two-frame safety layer over the bus with the inverted second
frame, then write down honestly what it does and does not give a node with no
assessment behind it.
\item Fill in the agency's worksheets completely for this node and publish them.
Complete, honest analyses of small systems are rare in public and cost nothing
but care.
\item Model this node's single channel in the free performance level calculator
and find out precisely where the architecture stops, rather than asserting it.
\item Build the second safe state for real: a brake, its release circuit, and a
deceleration that completes before power is removed. It is the half of this
chapter the bench could not demonstrate.
\end{itemize}

\subsection*{Roadmap and next steps}

Chapter 19 takes the node out of reach. Once a joint is inside a machine, the
only way to change its firmware is over the wire it already has, and a
bootloader that stays addressable is a different problem from one that runs on a
desk.

The published progression from here starts with the free material, because it is
better than its price suggests: the space agency's assurance handbook for
analysis method and worksheets, one silicon vendor's industrial functional safety
guide for the architecture arithmetic and the mechanism list, and the free
performance level calculator for modelling a structure. After that, the type-C
standard for robotics is the one that takes precedence for a joint, and the
emergency stop standard is short, cheap and the most directly useful of the paid
ones.

\subsection*{Portfolio evidence}

\begin{itemize}
\item The two safe states with the condition for each, which demonstrates
understanding the difference between an implementation and a risk assessment
output.
\item The document saying what the node is not.
\item The reachability and latch tests, which are short and make an unsourced
rule enforceable.
\item The filled-in failure analysis worksheets, from material that may be
published freely.
\item The evidence table, showing which claims were read from primary documents
and which were not.
\end{itemize}

\subsection*{Sources}

Normative references, with the correct current editions:
\begin{itemize}
\item IEC 61508-4:2010, edition 2.0, Friday 30 April 2010, "Functional safety of
electrical/electronic/programmable electronic safety-related systems, Part 4:
Definitions and abbreviations". Clause 3.1.13 for the safe state definition and
clause 3.2.11 for the tool classification, \textbf{both read from the official
preview}.
\item ISO 13850:2015, third edition, "Safety of machinery, Emergency stop
function, Principles for design", 11 pages. \textbf{Clause 4.1.1 read directly};
clause 4.1.3 on stop categories lies beyond the free preview and is cited by
clause number. \textbf{Its own normative references point at IEC 60204-1:2005},
which is two editions behind.
\item IEC 60204-1:2016 + AMD1:2021, edition 6.0, "Safety of machinery,
Electrical equipment of machines, Part 1: General requirements", clause 9.2.2
for the stop categories. \textbf{The category wording in this chapter is attested
by agreeing secondary sources, not primary verified.}
\item ISO 13849-1:2023, fourth edition, \pubdate{April 2023}, for performance levels. It
covers software design and \textbf{applies to high demand and continuous modes
only}.
\item IEC 62061:2021 + AMD1:2024, edition 2.0. \textbf{The 2021 edition did not
gain non-electrical coverage, it lost the restriction to electrical}: the older
title named electrical, electronic and programmable electronic systems and the
current one does not.
\item ISO 10218-1:2025 and ISO 10218-2:2025, \pubdate{February 2025}. \textbf{The series is
now titled "Robotics"}, not the older wording. For a joint this is the type-C
standard and \textbf{where a type-C standard differs, its provisions take
precedence}.
\item ISO/TS 15066:2016 on collaborative robots. \textbf{It is not withdrawn}:
published, confirmed in 2022, and current until its replacement appears.
\item IEC 61800-5-2:2016, edition 2.0, for the drive sub-functions. Edition 2
replaced safety function by \textbf{safety sub-function} throughout.
\textbf{Two brake-related sub-functions could not be confirmed on the page read,
and they are exactly what a gravity-loaded joint needs.}
\item IEC 60812:2018, edition 3.0, Friday 10 August 2018, for failure modes and
effects analysis. \textbf{Annex B.4 keeps the risk priority number} and adds an
alternative method beside it.
\item The coding guidelines: the \pubdate{March 2025} edition is current, the 2023 edition
consolidated the 2012 edition with its amendments and corrigendum, and the
compliance framework is mandatory from the 2023 edition onward.
\end{itemize}

Free and reusable material:
\begin{itemize}
\item The space agency's software engineering and assurance handbook, topic 8.5,
with four worksheet templates. \textbf{A United States government work with no
copyright notice: the safest material in this volume to adapt and publish.}
Attribute it.
\item The cancelled military standard that the method descends from, a United
States government work and quotable freely.
\item One silicon vendor's free industrial functional safety guide, for the
distinction between random hardware and systematic capability, the architecture
arithmetic, and the list of mechanisms implemented in firmware.
\item The free performance level calculator from an institute, version 3.0.4 of
Friday 31 October 2025. \textbf{Free to use and pass on; modification and
re-hosting are not permitted.}
\item A permissively licensed bus stack implementing the safety data object wire
format, Apache-2.0, which claims coding standard conformance with documented
exceptions and \textbf{makes no certification claim}.\newline
\url{https://github.com/CANopenNode/CANopenNode}
\item Stolte and co-authors, "A Taxonomy to Unify Fault Tolerance Regimes for
Automotive Systems", IEEE Transactions on Intelligent Vehicles volume 7 number
2, pages 251 to 262, 2022, which documents that the literature is inconsistent
and then repairs it. Cite and link; do not bundle.
\end{itemize}

Two gaps are recorded rather than glossed over. \textbf{The published safety
communication layer for this family of buses is written for the classic frame
format}, and this volume's research established no equivalent for the
flexible-data format. And the general framework for a safety layer over an
untrusted transport, in which the layer detects corruption, loss, repetition,
resequencing, delay and masquerade, is exactly the right frame for a host to
joint link; \textbf{the standard family's number, title and edition could not be
verified and are therefore not printed here.}
